% ============================================================
% S4: Theorem 4' -- Exact Constant Minimax Optimality
% Supplementary Information for Nature submission:
% "A Fundamental Impossibility in Data Quality"
%
% This is THE CROWN JEWEL: exact constant minimax optimality
% via Bahadur-Rao asymptotics, Chernoff information, and the
% adaptive threshold achieving the lower bound constant.
% ============================================================

\section{Theorem~4': Exact Constant Minimax Optimality}
\label{sec:thm4}

\subsection{Statement of the Theorem}
\label{sec:thm4_statement}

Theorem~1 establishes that the SCX noise detector achieves an exponential
convergence rate $2M\Delta^2$.  Theorem~4' refines this result in three
ways: it replaces the Hoeffding rate $2\Delta^2$ with the exact Chernoff
information $\kappa$, it establishes the \emph{exact constant} multiplying the
exponential, and it proves that no algorithm can improve upon this constant.

Throughout this section we work within a single state $s$ satisfying
Assumptions~(A1)--(A6).  Let
\[
p_0 = \mu_s,\qquad
p_1 = 1 - C_{\text{bal}}\cdot\frac{\mu_s}{K-1},
\]
where $K=K_{\mathcal{Y}}$ is the number of classes and
$0<p_0<p_1<1$.  The expert error indicators
$e_m=\mathbf{1}\{f_m(x)\neq y\}$ are i.i.d.\ $\operatorname{Bern}(p_0)$
under the clean hypothesis $\mathrm{H}_0$ and i.i.d.\ $\operatorname{Bern}(p_1)$
under the noise hypothesis $\mathrm{H}_1$.

Define the following quantities:
\[
\begin{aligned}
\kappa &= C\bigl(\operatorname{Bern}(p_0),\operatorname{Bern}(p_1)\bigr)
        &&\text{Chernoff information},\\[4pt]
\theta^* &= \frac{\log\frac{1-p_0}{1-p_1}}
                 {\log\frac{p_1(1-p_0)}{p_0(1-p_1)}}
        &&\text{Chernoff point (unique root of
          $\operatorname{KL}(\theta\|p_0)=\operatorname{KL}(\theta\|p_1)$)},\\[4pt]
\lambda_0^* &= \log\frac{\theta^*(1-p_0)}{p_0(1-\theta^*)}>0,
\qquad
\lambda_1^* = \log\frac{\theta^*(1-p_1)}{p_1(1-\theta^*)}<0
        &&\text{Saddlepoints at $\theta^*$},\\[4pt]
D^* &= \lambda_0^* + |\lambda_1^*|
      = \log\frac{p_1(1-p_0)}{p_0(1-p_1)}>0,\\[4pt]
s &= \frac{|\lambda_1^*|}{D^*}\in(0,1).
\end{aligned}
\]

\begin{theorem}[Exact Constant Minimax Optimality of SCX]
  \label{thm:exact_constant}
  Under Assumptions~(A1)--(A6), for any state $s$ with $p_0<p_1$, the
  following hold.
  \begin{enumerate}[(a)]
    \item \textbf{SCX achievability with adaptive threshold.}
      The SCX detector using threshold
      \[
      \theta^\dagger = \theta^*
        + \frac{1}{M}\frac{\log\frac{1-\eta}{\eta}}{D^*}
        + O\!\left(\frac1{M^2}\right)
      \]
      achieves
      \[
      \lim_{M\to\infty}
        e^{M\kappa}\sqrt{2\pi M}\,
        \bigl(1-\operatorname{F1}_{\text{SCX}}(\theta^\dagger)\bigr)
      = \frac{C_{\min}}{\eta},
      \]
      where the canonical minimax constant is
      \[
      \boxed{\;
      C_{\min}
      = \frac{\eta}{2}
        \left(\frac{1-\eta}{\eta}\right)^{\!s}
        \frac{1/\lambda_0^* + 1/|\lambda_1^*|}
             {\sqrt{\theta^*(1-\theta^*)}}\;}.
      \]

    \item \textbf{Minimax lower bound.}
      For \emph{any} noise detection algorithm $\mathcal{A}$,
      \[
      \liminf_{M\to\infty}
        e^{M\kappa}\sqrt{2\pi M}\,
        \bigl(1-\operatorname{F1}_{\mathcal{A}}\bigr)
      \ge \frac{C_{\min}}{\eta}.
      \]

    \item \textbf{Constant optimality.}
      Since SCX with the adaptive threshold $\theta^\dagger$ attains the
      limit with constant $C_{\min}$, it is \emph{exact constant minimax
      optimal}: no algorithm can achieve a smaller asymptotic error constant.

    \item \textbf{Multi-state generalisation.}
      For $S$ states with proportions $\rho_s$, Chernoff information
      $\kappa_s$, and per-state constants $C_s$,
      \[
      \kappa_{\text{global}} = \min_{1\le s\le S}\kappa_s,\qquad
      C_{\text{global}} = \sum_{s:\,\kappa_s=\kappa_{\text{global}}} \rho_s C_s,
      \]
      and $\displaystyle\liminf_{M\to\infty}
        e^{M\kappa_{\text{global}}}\sqrt{2\pi M}\,
        (1-\operatorname{F1}_{\mathcal{A}})
        \ge C_{\text{global}}/\eta$.
  \end{enumerate}
\end{theorem}

The remainder of this section provides the complete proof, organised into
lemmas that are of independent interest.

\subsection{Preliminaries: Chernoff Information and Bahadur-Rao}
\label{sec:thm4_prelim}

\subsubsection{Chernoff Information}

For two distributions $P_0,P_1$ with densities $p_0(\cdot),p_1(\cdot)$,
the Chernoff information is
\[
C(P_0,P_1) = -\min_{\lambda\in[0,1]}
  \log\mathbb{E}_{P_0}\!\left[\left(\frac{dP_1}{dP_0}\right)^{\!\!\lambda}\right].
\]
For $\operatorname{Bern}(p_0)$ vs.\ $\operatorname{Bern}(p_1)$, we have
\[
\mathbb{E}_{P_0}\!\left[\left(\frac{dP_1}{dP_0}\right)^{\!\!\lambda}\right]
= p_0^{1-\lambda}p_1^\lambda + (1-p_0)^{1-\lambda}(1-p_1)^\lambda.
\]

\subsubsection{Bahadur-Rao Theorem (Bernoulli Case)}

The Bahadur-Rao theorem \cite{bahadur1960deviations} provides the exact
asymptotic probability of a large deviation, including the constant
prefactor.  For i.i.d.\ $\operatorname{Bern}(p)$ variables
$X_1,\dots,X_M$ and a threshold $\theta>p$,
\[
\mathbb{P}\!\left(\frac1M\sum_{i=1}^M X_i \ge \theta\right)
\sim
\frac{\exp\!\bigl(-M\cdot\operatorname{KL}(\theta\|p)\bigr)}
     {\lambda^*(\theta)\,\sqrt{2\pi M\,\theta(1-\theta)}},
\qquad M\to\infty,
\]
where $\lambda^*(\theta)=\log\frac{\theta(1-p)}{p(1-\theta)}>0$ is the
Cram\'er saddlepoint.  For the lower tail $\theta<p$, the same formula
holds with $|\lambda^*(\theta)|$ in the denominator.

\subsubsection{Exponential Tilting}

The proof of the Bahadur-Rao theorem uses exponential tilting.  Define the
tilted measure $\mathbb{P}_{\lambda^*}$ by
\[
\frac{d\mathbb{P}_{\lambda^*}}{d\mathbb{P}}
= \exp\!\bigl(\lambda^* S_M - M\psi(\lambda^*)\bigr),\qquad
S_M=\sum_{i=1}^M X_i,
\]
where $\psi(\lambda)=\log(1-p+pe^\lambda)$.  Under $\mathbb{P}_{\lambda^*}$,
the $X_i$ are i.i.d.\ $\operatorname{Bern}(\theta)$, i.e.\ the tilted
distribution has mean exactly $\theta$.  The saddlepoint $\lambda^*$ and
variance $\sigma^2=\theta(1-\theta)$ are obtained from
\[
\psi'(\lambda^*)=\theta,\qquad
\psi''(\lambda^*)=\theta(1-\theta).
\]

\subsection{Part I: Achievability (SCX Upper Bound)}
\label{sec:thm4_upper}

\subsubsection{Lemma~A: Bahadur-Rao Exact Asymptotics for Bernoulli}
\label{sec:lemmaA}

\begin{lemma}[Bahadur-Rao for Bernoulli, exact form]
  \label{lem:BahadurRao}
  Let $X_1,\dots,X_M\stackrel{\text{i.i.d.}}{\sim}\operatorname{Bern}(p)$ and
  fix $\theta>p$.  Then
  \[
  \mathbb{P}\!\left(\frac1M\sum_{i=1}^M X_i \ge \theta\right)
  = \frac{\exp\!\bigl(-M\cdot\operatorname{KL}(\theta\|p)\bigr)}
         {\lambda^*\sqrt{2\pi M\,\theta(1-\theta)}}
    \bigl(1+O(M^{-1})\bigr),
  \]
  where $\lambda^*=\log\frac{\theta(1-p)}{p(1-\theta)}>0$ and the $O(M^{-1})$
  constant is explicit in $p$ and $\theta$.
\end{lemma}

\begin{proof}[Sketch]
  Write the binomial tail probability
  $\mathbb{P}(S_M\ge k)=\sum_{j=k}^M\binom{M}{j}p^j(1-p)^{M-j}$ with
  $k=\lceil M\theta\rceil$.  Apply Stirling's formula with Robbins' remainder
  bound to the leading term $a_k$, then bound the ratio of consecutive terms
  $a_{j+1}/a_j\le e^{-\lambda^*}(1+O(M^{-1}))$ and sum the geometric series.
  The lattice correction factor $(1-e^{-\lambda^*})^{-1}$ is absorbed into
  the $O(M^{-1})$ term; the form $1/\lambda^*$ is obtained by
  $\lambda^*/(1-e^{-\lambda^*})=1+O(\lambda^*)$.
\end{proof}

For the lower tail $\theta<p$, symmetry gives
\[
\mathbb{P}\!\left(\frac1M\sum_{i=1}^M X_i \le \theta\right)
= \frac{\exp\!\bigl(-M\cdot\operatorname{KL}(\theta\|p)\bigr)}
       {|\lambda^*|\sqrt{2\pi M\,\theta(1-\theta)}}
  \bigl(1+O(M^{-1})\bigr),
\]
where $|\lambda^*|=\log\frac{p(1-\theta)}{\theta(1-p)}$.

\subsubsection{Application to FPR and FNR}

Apply Lemma~A to the expert error indicators.  Under $\mathrm{H}_0$ (clean),
$e_m\sim\operatorname{Bern}(p_0)$.  Under $\mathrm{H}_1$ (noise),
$e_m\sim\operatorname{Bern}(p_1)$.  For a fixed threshold
$\theta\in(p_0,p_1)$,
\[
\begin{aligned}
\operatorname{FPR}_M(\theta)
&= \mathbb{P}(C_M\ge\theta\mid\text{clean})
   \sim \frac{\exp\!\bigl(-M\cdot\operatorname{KL}(\theta\|p_0)\bigr)}
            {\lambda_0^*\sqrt{2\pi M\,\theta(1-\theta)}},\\[4pt]
\operatorname{FNR}_M(\theta)
&= \mathbb{P}(C_M\le\theta\mid\text{noise})
   \sim \frac{\exp\!\bigl(-M\cdot\operatorname{KL}(\theta\|p_1)\bigr)}
            {|\lambda_1^*|\sqrt{2\pi M\,\theta(1-\theta)}},
\end{aligned}
\]
where $\lambda_0^*=\log\frac{\theta(1-p_0)}{p_0(1-\theta)}>0$ and
$\lambda_1^*=\log\frac{\theta(1-p_1)}{p_1(1-\theta)}<0$,
$|\lambda_1^*|=\log\frac{p_1(1-\theta)}{\theta(1-p_1)}$.

\subsubsection{Lemma~B: F1 Asymptotic Expansion}
\label{sec:lemmaB}

\begin{lemma}[F1 expansion]
  \label{lem:F1expansion}
  Let $\eta=\mathbb{P}(\text{noise})$, $\operatorname{FPR}=\operatorname{FPR}_M$,
  $\operatorname{FNR}=\operatorname{FNR}_M$, and define
  $r=\operatorname{FNR}/2+(1-\eta)\operatorname{FPR}/(2\eta)$.  For
  $M$ sufficiently large that $r\le\frac12$,
  \[
  1-\operatorname{F1}
  = \frac{\operatorname{FNR}}{2}
    + \frac{1-\eta}{2\eta}\operatorname{FPR}
    + R,
  \]
  where $|R|\le 2r^2$.
\end{lemma}

\begin{proof}
  From the definitions,
  \[
  \operatorname{F1}
  = \frac{2\eta\cdot\operatorname{TPR}}
         {\eta(1+\operatorname{TPR})+(1-\eta)\operatorname{FPR}}
  = \frac{2\eta(1-\operatorname{FNR})}
         {\eta(2-\operatorname{FNR})+(1-\eta)\operatorname{FPR}}.
  \]
  Direct algebra gives
  \[
  1-\operatorname{F1}
  = \frac{\eta\operatorname{FNR}+(1-\eta)\operatorname{FPR}}
         {\eta(2-\operatorname{FNR})+(1-\eta)\operatorname{FPR}}.
  \]
  Let $\alpha=\eta\operatorname{FNR}+(1-\eta)\operatorname{FPR}$ and
  $\beta=2\eta-\eta\operatorname{FNR}+(1-\eta)\operatorname{FPR}$.  Since
  $\beta\ge\eta>0$, write
  $1-\operatorname{F1} = \frac{\alpha}{2\eta}\cdot\frac1{1-\gamma}$
  with $\gamma=\frac{\operatorname{FNR}}2-\frac{(1-\eta)\operatorname{FPR}}{2\eta}$.
  For $|\gamma|<1$ (which holds when $r\le\frac12$), expand the geometric
  series to obtain the first-order terms and bound the remainder.
\end{proof}

Substituting the Bahadur-Rao asymptotics at a threshold $\theta$,
\[
1-\operatorname{F1}(\theta)
= \frac1{\sqrt{2\pi M\,\theta(1-\theta)}}
  \Bigl[
    \frac{e^{-M\cdot\operatorname{KL}(\theta\|p_1)}}{2|\lambda_1^*(\theta)|}
    + \frac{1-\eta}{2\eta}\,
      \frac{e^{-M\cdot\operatorname{KL}(\theta\|p_0)}}{\lambda_0^*(\theta)}
  \Bigr]
  + O\!\left(\frac{e^{-2M\min(\kappa_0,\kappa_1)}}{M}\right),
\]
where $\kappa_0=\operatorname{KL}(\theta\|p_0)$ and
$\kappa_1=\operatorname{KL}(\theta\|p_1)$.  The remainder decays twice as
fast exponentially and is asymptotically negligible.

\subsubsection{Lemma~C: Chernoff Information Closed Form}
\label{sec:lemmaC}

\begin{lemma}[Chernoff point and information]
  \label{lem:Chernoff}
  The unique $\theta^*\in(p_0,p_1)$ satisfying
  $\operatorname{KL}(\theta^*\|p_0)=\operatorname{KL}(\theta^*\|p_1)$ is
  \[
  \boxed{\;
  \theta^*
  = \frac{\log\frac{1-p_0}{1-p_1}}
         {\log\frac{p_1(1-p_0)}{p_0(1-p_1)}}\;}.
  \]
  The common value $\kappa=\operatorname{KL}(\theta^*\|p_0)$ is the
  Chernoff information $C(\operatorname{Bern}(p_0),\operatorname{Bern}(p_1))$.
\end{lemma}

\begin{proof}
  Equating the KL divergences and cancelling the common term
  $\theta\log\theta+(1-\theta)\log(1-\theta)$ gives
  \[
  \theta\log\frac1{p_0}+(1-\theta)\log\frac1{1-p_0}
  = \theta\log\frac1{p_1}+(1-\theta)\log\frac1{1-p_1}.
  \]
  Simplifying:
  \[
  \theta\log\frac{p_1}{p_0} = (1-\theta)\log\frac{1-p_0}{1-p_1},
  \]
  which yields the closed form.  Strict convexity of $\operatorname{KL}$
  in its first argument and the sign change
  $f(p_0)=-\operatorname{KL}(p_0\|p_1)<0$,
  $f(p_1)=\operatorname{KL}(p_1\|p_0)>0$ guarantee uniqueness and
  $\theta^*\in(p_0,p_1)$.
\end{proof}

The saddlepoints at $\theta^*$ are:
\[
\lambda_0^* = \log\frac{\theta^*(1-p_0)}{p_0(1-\theta^*)}>0,\qquad
\lambda_1^* = \log\frac{\theta^*(1-p_1)}{p_1(1-\theta^*)}<0.
\]

\subsubsection{Lemma~D: Adaptive Threshold Optimality}
\label{sec:lemmaD}

The naive threshold $\theta^*$ (the Chernoff point) balances the two
exponential rates but ignores the noise prior $\eta$.  When $\eta\neq\frac12$,
the optimal threshold---the one minimising $1-\operatorname{F1}$---differs
from $\theta^*$ by $O(1/M)$.  This $O(1/M)$ shift produces an $O(1)$
multiplicative factor in the error probability via the exponent, which is
essential for achieving the minimax lower bound.

\begin{lemma}[Adaptive threshold]
  \label{lem:adaptive_threshold}
  The threshold $\theta^\dagger$ that minimises the leading-order expression
  for $1-\operatorname{F1}$ satisfies
  \[
  \operatorname{KL}(\theta^\dagger\|p_0)
  = \operatorname{KL}(\theta^\dagger\|p_1)
    + \frac1M\log\frac{1-\eta}{\eta}
    + O\!\left(\frac1{M^2}\right).
  \]
  Expanding around $\theta^*$ gives
  \[
  \theta^\dagger = \theta^*
    + \frac1M\,\frac{\log\frac{1-\eta}{\eta}}{D^*}
    + O\!\left(\frac1{M^2}\right),
  \]
  where $D^*=\lambda_0^*+|\lambda_1^*|
          =\log\frac{p_1(1-p_0)}{p_0(1-p_1)}$.
\end{lemma}

\begin{proof}
  The objective function (ignoring the common factor
  $1/\sqrt{2\pi M\theta(1-\theta)}$) is
  \[
  \Phi_M(\theta)
  = \frac{e^{-M\cdot\operatorname{KL}(\theta\|p_1)}}{2|\lambda_1^*(\theta)|}
    + \frac{1-\eta}{2\eta}\,
      \frac{e^{-M\cdot\operatorname{KL}(\theta\|p_0)}}{\lambda_0^*(\theta)}.
  \]
  Differentiating and keeping only the dominant
  $-M\operatorname{KL}'(\theta\|p)=-M\lambda^*(\theta)$ terms (the
  $-M$ factor dominates the $O(1)$ derivative of $\log\lambda^*$),
  the first-order condition reduces to
  \[
  e^{-M\cdot\operatorname{KL}(\theta\|p_1)}
  \approx \frac{1-\eta}{\eta}\,
          e^{-M\cdot\operatorname{KL}(\theta\|p_0)}.
  \]
  Taking logs yields the stated relation.  Expanding
  $\operatorname{KL}(\theta^*+\delta\|p_0)-\operatorname{KL}(\theta^*+\delta\|p_1)$
  gives $\delta D^* + O(\delta^3)$, where the quadratic term cancels exactly
  because $\operatorname{KL}''(\cdot\|p_0)=\operatorname{KL}''(\cdot\|p_1)
  =1/(\theta(1-\theta))$ is independent of $p$.  Solving $\delta D^* =
  \frac1M\log\frac{1-\eta}{\eta}+O(M^{-2})$ gives the result.
\end{proof}

\paragraph{The $O(1)$ exponential prefactor.}
The $O(1/M)$ shift in $\theta$ produces an $O(1)$ multiplicative factor in
$\exp(-M\cdot\operatorname{KL})$.  Expanding the KL divergence at
$\theta^\dagger=\theta^*+\delta$:
\[
\begin{aligned}
\operatorname{KL}(\theta^\dagger\|p_0)
&= \kappa + \lambda_0^*\delta
   + \frac{\delta^2}{2\theta^*(1-\theta^*)} + O(\delta^3),\\[4pt]
\operatorname{KL}(\theta^\dagger\|p_1)
&= \kappa - |\lambda_1^*|\delta
   + \frac{\delta^2}{2\theta^*(1-\theta^*)} + O(\delta^3).
\end{aligned}
\]

Multiplying by $M$ and substituting $\delta$:
\[
\begin{aligned}
M\cdot\operatorname{KL}(\theta^\dagger\|p_0)
&= M\kappa + \frac{\lambda_0^*}{D^*}\log\frac{1-\eta}{\eta} + O\!\left(\frac1M\right),\\[4pt]
M\cdot\operatorname{KL}(\theta^\dagger\|p_1)
&= M\kappa - \frac{|\lambda_1^*|}{D^*}\log\frac{1-\eta}{\eta} + O\!\left(\frac1M\right).
\end{aligned}
\]

Exponentiating:
\[
\begin{aligned}
\exp\!\bigl(-M\cdot\operatorname{KL}(\theta^\dagger\|p_0)\bigr)
&= e^{-M\kappa}
   \left(\frac{1-\eta}{\eta}\right)^{\!-\lambda_0^*/D^*}
   \bigl(1+O(M^{-1})\bigr),\\[4pt]
\exp\!\bigl(-M\cdot\operatorname{KL}(\theta^\dagger\|p_1)\bigr)
&= e^{-M\kappa}
   \left(\frac{1-\eta}{\eta}\right)^{\!|\lambda_1^*|/D^*}
   \bigl(1+O(M^{-1})\bigr).
\end{aligned}
\]

\paragraph{FPR and FNR at the adaptive threshold.}

Substituting into the Bahadur-Rao expansions:
\[
\begin{aligned}
\operatorname{FPR}_M(\theta^\dagger)
&\sim \frac{e^{-M\kappa}}
          {\lambda_0^*\sqrt{2\pi M\,\theta^*(1-\theta^*)}}\,
        \left(\frac{1-\eta}{\eta}\right)^{\!-\lambda_0^*/D^*},\\[4pt]
\operatorname{FNR}_M(\theta^\dagger)
&\sim \frac{e^{-M\kappa}}
          {|\lambda_1^*|\sqrt{2\pi M\,\theta^*(1-\theta^*)}}\,
        \left(\frac{1-\eta}{\eta}\right)^{\!|\lambda_1^*|/D^*}.
\end{aligned}
\]

\paragraph{The critical cancellation in $1-\operatorname{F1}$.}

Recall $s=|\lambda_1^*|/D^*$ and $\lambda_0^*/D^*=1-s$.  The FNR contribution to
$1-\operatorname{F1}$ is
\[
\frac12\operatorname{FNR}_M(\theta^\dagger)
\sim \frac{e^{-M\kappa}}
          {2|\lambda_1^*|\sqrt{2\pi M\,\theta^*(1-\theta^*)}}\,
      \left(\frac{1-\eta}{\eta}\right)^{\!s}.
\]

The FPR contribution is
\[
\begin{aligned}
\frac{1-\eta}{2\eta}\operatorname{FPR}_M(\theta^\dagger)
&\sim \frac{1-\eta}{2\eta}\,
      \frac{e^{-M\kappa}}
           {\lambda_0^*\sqrt{2\pi M\,\theta^*(1-\theta^*)}}\,
      \left(\frac{1-\eta}{\eta}\right)^{\!-(1-s)}\\[4pt]
&= \frac{e^{-M\kappa}}
        {2\lambda_0^*\sqrt{2\pi M\,\theta^*(1-\theta^*)}}\,
      \left(\frac{1-\eta}{\eta}\right)^{\!s}.
\end{aligned}
\]

\textbf{Both terms carry the identical factor
$\bigl((1-\eta)/\eta\bigr)^s$!}  This is the critical cancellation that
makes the adaptive threshold constant-optimal.

\paragraph{Asymptotic constant for SCX.}

Summing the two contributions:
\[
\begin{aligned}
1-\operatorname{F1}_{\text{SCX}}(\theta^\dagger)
&\sim \frac{e^{-M\kappa}}
          {\sqrt{2\pi M\,\theta^*(1-\theta^*)}}\,
      \left(\frac{1-\eta}{\eta}\right)^{\!s}
      \frac12\!\left(
        \frac1{\lambda_0^*} + \frac1{|\lambda_1^*|}
      \right).
\end{aligned}
\]

Multiplying by $e^{M\kappa}\sqrt{2\pi M}$ and taking the limit:
\[
\boxed{\;
\lim_{M\to\infty}
  e^{M\kappa}\sqrt{2\pi M}\,
  \bigl(1-\operatorname{F1}_{\text{SCX}}(\theta^\dagger)\bigr)
= \frac12\left(\frac{1-\eta}{\eta}\right)^{\!s}
  \frac{1/\lambda_0^* + 1/|\lambda_1^*|}
       {\sqrt{\theta^*(1-\theta^*)}}\;}.
\]

This equals $C_{\min}/\eta$ because
\[
\frac{C_{\min}}{\eta}
= \frac12\left(\frac{1-\eta}{\eta}\right)^{\!s}
  \frac{1/\lambda_0^* + 1/|\lambda_1^*|}
       {\sqrt{\theta^*(1-\theta^*)}}.
\]

Thus Lemma~D establishes part~(a) of Theorem~4'.

\subsection{Part II: Lower Bound (Minimax Converse)}
\label{sec:thm4_lower}

\subsubsection{Lemma~E: Second-Order Asymptotic Lower Bound}
\label{sec:lemmaE}

\begin{lemma}[Exact constant minimax lower bound]
  \label{lem:lowerbound}
  For any noise detection algorithm $\mathcal{A}$,
  \[
  \liminf_{M\to\infty}
    e^{M\kappa}\sqrt{2\pi M}\,
    \bigl(1-\operatorname{F1}_{\mathcal{A}}(M)\bigr)
  \ge \frac{C_{\min}}{\eta},
  \]
  where $C_{\min}$ is the constant defined in Theorem~4'.
  Equivalently, writing $r=\lambda_0^*/D^*=1-s$,
  \[
  C_{\min}
  = \frac{\eta^r(1-\eta)^{1-r}}2\,
    \frac{1/\lambda_0^*+1/|\lambda_1^*|}
         {\sqrt{\theta^*(1-\theta^*)}}.
  \]
\end{lemma}

\begin{proof}
  The proof proceeds in five parts.

  \paragraph{Part 1: Reduction to a binary hypothesis test.}
  Within a fixed state $s$, the expert errors $e_1,\dots,e_M$ are i.i.d.\
  $\operatorname{Bern}(p_0)$ under $\mathrm{H}_0$ (clean) and
  $\operatorname{Bern}(p_1)$ under $\mathrm{H}_1$ (noise).  Any algorithm
  $\mathcal{A}$ reduces to a decision rule
  $\delta_M:\{0,1\}^M\to\{0,1\}$ with Type~I error
  $\alpha_M=\Pr(\delta_M=1\mid\mathrm{H}_0)$ and Type~II error
  $\beta_M=\Pr(\delta_M=0\mid\mathrm{H}_1)$.

  From Lemma~B,
  \[
  1-\operatorname{F1}_{\mathcal{A}}
  = \frac{1-\eta}{2\eta}\,\alpha_M + \frac12\,\beta_M + o(1),
  \]
  where the $o(1)$ term is uniform over decision rules and decays as
  $O(e^{-2M\kappa})$, negligible at the leading order.  Define the
  weighted risk
  \[
  R_M(\delta_M) = w_0\alpha_M + w_1\beta_M,\qquad
  w_0=\frac{1-\eta}{2\eta},\; w_1=\frac12.
  \]

  \paragraph{Part 2: Neyman-Pearson reduction.}
  By the Neyman-Pearson lemma, for any test with Type~I error $\le\alpha$,
  the minimum achievable Type~II error is attained by the likelihood ratio
  test (LRT).  Therefore for any decision rule $\delta_M$,
  \[
  R_M(\delta_M) \ge \min_{\alpha\in[0,1]}
    \bigl\{w_0\alpha + w_1\beta_{\text{NP}}(\alpha)\bigr\},
  \]
  where $\beta_{\text{NP}}(\alpha)$ is the Type~II error of the most powerful
  level-$\alpha$ test.

  The minimiser $\alpha^*$ satisfies $w_0 + w_1\beta_{\text{NP}}'(\alpha^*)=0$.
  For the Neyman-Pearson test, the slope of the ROC curve at $\alpha^*$ is
  $-1/\tau^*$, where $\tau^*$ is the likelihood ratio threshold.  Hence
  $\tau^* = w_0/w_1 = (1-\eta)/\eta$; the minimiser is exactly the Bayes
  decision rule with prior odds $(1-\eta):\eta$ against $\mathrm{H}_0$.  Thus
  for any $\mathcal{A}$,
  \[
  R_M(\delta_{\mathcal{A}})
  \ge R_M^* := w_0\alpha_M^* + w_1\beta_M^*,
  \]
  where $\delta_M^*(e)=\mathbf{1}\{\log L_M(e) > \log\frac{1-\eta}{\eta}\}$
  and $L_M(e)=\prod_{m=1}^M\frac{p_1^{e_m}(1-p_1)^{1-e_m}}{p_0^{e_m}(1-p_0)^{1-e_m}}$.

  \paragraph{Part 3: Exact asymptotics of the Bayes test.}
  The log-likelihood ratio simplifies to
  \[
  \log L_M = M\bigl(C_M\Delta + \Delta_0\bigr),
  \]
  where
  $\Delta=\log\frac{p_1(1-p_0)}{p_0(1-p_1)} = D^*$ and
  $\Delta_0=\log\frac{1-p_1}{1-p_0}$.
  The rejection region $C_M > -\Delta_0/\Delta + \frac1{M\Delta}\log\frac{1-\eta}{\eta}$
  has threshold
  \[
  t_M = \theta^* + \frac1{M D^*}\log\frac{1-\eta}{\eta}.
  \]
  This is precisely the adaptive threshold $\theta^\dagger$ from Lemma~D.

  The error probabilities of the Bayes test are
  \[
  \alpha_M^* = \Pr\nolimits_{p_0}(C_M > t_M),\qquad
  \beta_M^*  = \Pr\nolimits_{p_1}(C_M \le t_M),
  \]
  up to a $O(e^{-M\kappa}/M)$ correction for tie-breaking.

  Apply the Bahadur-Rao theorem (Lemma~A) to each:
  \[
  \begin{aligned}
  \alpha_M^* &\sim
    \frac{\exp\!\bigl(-M\operatorname{KL}(t_M\|p_0)\bigr)}
         {\lambda_0^*(t_M)\sqrt{2\pi M\,t_M(1-t_M)}},\\[4pt]
  \beta_M^*  &\sim
    \frac{\exp\!\bigl(-M\operatorname{KL}(t_M\|p_1)\bigr)}
         {|\lambda_1^*(t_M)|\sqrt{2\pi M\,t_M(1-t_M)}}.
  \end{aligned}
  \]

  \paragraph{Part 4: Expansion around the Chernoff point.}
  Write $\delta = t_M-\theta^* = \frac1{M D^*}\log\frac{1-\eta}{\eta} + O(M^{-2})$.
  Expand the KL divergences as in Lemma~D:
  \[
  \begin{aligned}
  M\operatorname{KL}(t_M\|p_0)
  &= M\kappa + \frac{\lambda_0^*}{D^*}\log\frac{1-\eta}{\eta} + o(1),\\[4pt]
  M\operatorname{KL}(t_M\|p_1)
  &= M\kappa - \frac{|\lambda_1^*|}{D^*}\log\frac{1-\eta}{\eta} + o(1).
  \end{aligned}
  \]

  Substituting into the Bahadur-Rao expressions:
  \[
  \begin{aligned}
  \alpha_M^* &=
    \frac{e^{-M\kappa}}
         {\lambda_0^*\sqrt{2\pi M\,\theta^*(1-\theta^*)}}\,
    \left(\frac{1-\eta}{\eta}\right)^{\!-\lambda_0^*/D^*}
    \bigl(1+o(1)\bigr),\\[4pt]
  \beta_M^* &=
    \frac{e^{-M\kappa}}
         {|\lambda_1^*|\sqrt{2\pi M\,\theta^*(1-\theta^*)}}\,
    \left(\frac{1-\eta}{\eta}\right)^{\!|\lambda_1^*|/D^*}
    \bigl(1+o(1)\bigr).
  \end{aligned}
  \]

  The weighted risk is therefore
  \[
  \begin{aligned}
  R_M^* &=
    w_0\alpha_M^* + w_1\beta_M^* \\[4pt]
  &= \frac{e^{-M\kappa}}
          {\sqrt{2\pi M\,\theta^*(1-\theta^*)}}\,
     \Bigg[
       \frac{1-\eta}{2\eta}\,
       \frac1{\lambda_0^*}
       \left(\frac{1-\eta}{\eta}\right)^{\!-\lambda_0^*/D^*}
     + \frac12\,
       \frac1{|\lambda_1^*|}
       \left(\frac{1-\eta}{\eta}\right)^{\!|\lambda_1^*|/D^*}
     \Bigg] \bigl(1+o(1)\bigr).
  \end{aligned}
  \]

  Using $s=|\lambda_1^*|/D^*$ and $\lambda_0^*/D^*=1-s$, the terms simplify:
  \[
  \begin{aligned}
  \frac{1-\eta}{2\eta\lambda_0^*}
    \left(\frac{1-\eta}{\eta}\right)^{\!-(1-s)}
  &= \frac1{2\lambda_0^*}
     \left(\frac{1-\eta}{\eta}\right)^{\!s},\\[4pt]
  \frac1{2|\lambda_1^*|}
    \left(\frac{1-\eta}{\eta}\right)^{\!s}
  &= \frac1{2|\lambda_1^*|}
     \left(\frac{1-\eta}{\eta}\right)^{\!s}.
  \end{aligned}
  \]

  Both terms share the common factor
  $\frac12\bigl((1-\eta)/\eta\bigr)^s$.  Factoring it out:
  \[
  R_M^* = \frac{e^{-M\kappa}}
              {\sqrt{2\pi M\,\theta^*(1-\theta^*)}}\,
          \frac12\left(\frac{1-\eta}{\eta}\right)^{\!s}
          \left(\frac1{\lambda_0^*} + \frac1{|\lambda_1^*|}\right)
          \bigl(1+o(1)\bigr).
  \]

  \paragraph{Part 5: Conversion to the F1 lower bound.}
  From Part~1, $1-\operatorname{F1}_{\mathcal{A}} = R_M(\delta_{\mathcal{A}}) + o(1)
  \ge R_M^* + o(1)$.  Multiplying by $e^{M\kappa}\sqrt{2\pi M}$ and taking the
  limit inferior:
  \[
  \liminf_{M\to\infty}
    e^{M\kappa}\sqrt{2\pi M}\,
    \bigl(1-\operatorname{F1}_{\mathcal{A}}(M)\bigr)
  \ge \frac12\left(\frac{1-\eta}{\eta}\right)^{\!s}
      \frac{1/\lambda_0^* + 1/|\lambda_1^*|}
           {\sqrt{\theta^*(1-\theta^*)}}.
  \]

  The right-hand side equals $C_{\min}/\eta$ by definition of $C_{\min}$.
  This completes the proof of Lemma~E. \hfill$\square$
\end{proof}

\subsection{Part III: Constant Matching $\to$ Optimality}
\label{sec:thm4_matching}

\subsubsection{The $O(1)$ Prefactor Cancellation}

The central mechanism of the constant optimality proof is the cancellation
that produces the identical factor $((1-\eta)/\eta)^s$ in both the FPR and
FNR contributions to $1-\operatorname{F1}$.

The adaptive threshold $\theta^\dagger$ differs from the Chernoff point
$\theta^*$ by $O(1/M)$.  This $O(1/M)$ shift generates an $O(1)$ change in
the exponent:
\[
\exp\!\bigl(-M\cdot\operatorname{KL}(\theta^\dagger\|p)\bigr)
= e^{-M\kappa} \cdot \text{constant factor}.
\]

The constant factor for $p=p_0$ is $((1-\eta)/\eta)^{-\lambda_0^*/D^*}$
and for $p=p_1$ is $((1-\eta)/\eta)^{|\lambda_1^*|/D^*}$.  When combined
with the F1 weights $w_0=(1-\eta)/(2\eta)$ and $w_1=1/2$, both contributions
acquire the identical factor $((1-\eta)/\eta)^s$ because:
\[
\frac{1-\eta}{2\eta}
\cdot \left(\frac{1-\eta}{\eta}\right)^{\!-\lambda_0^*/D^*}
= \frac1{2\lambda_0^*}
  \left(\frac{1-\eta}{\eta}\right)^{\!s},
\qquad
\frac12
\cdot \left(\frac{1-\eta}{\eta}\right)^{\!|\lambda_1^*|/D^*}
= \frac1{2|\lambda_1^*|}
  \left(\frac{1-\eta}{\eta}\right)^{\!s}.
\]

This cancellation is mathematically unavoidable---it follows from the
relationship $s+(1-s)=1$, i.e.,
$|\lambda_1^*|/D^* + \lambda_0^*/D^* = 1$, which is a consequence of the
definition $D^*=\lambda_0^*+|\lambda_1^*|$.

\subsubsection{Proof that SCX Achieves $C_{\min}$}

From Lemma~D, the SCX detector with adaptive threshold $\theta^\dagger$
attains:
\[
\lim_{M\to\infty}
  e^{M\kappa}\sqrt{2\pi M}\,
  \bigl(1-\operatorname{F1}_{\text{SCX}}(\theta^\dagger)\bigr)
= \frac12\left(\frac{1-\eta}{\eta}\right)^{\!s}
  \frac{1/\lambda_0^* + 1/|\lambda_1^*|}
       {\sqrt{\theta^*(1-\theta^*)}}.
\]

From Lemma~E, the minimax lower bound for any algorithm is:
\[
\liminf_{M\to\infty}
  e^{M\kappa}\sqrt{2\pi M}\,
  \bigl(1-\operatorname{F1}_{\mathcal{A}}\bigr)
\ge \frac12\left(\frac{1-\eta}{\eta}\right)^{\!s}
    \frac{1/\lambda_0^* + 1/|\lambda_1^*|}
         {\sqrt{\theta^*(1-\theta^*)}}.
\]

The two expressions are \textbf{identical}.  Therefore the SCX detector
attains the lower bound, establishing exact constant minimax optimality.
This proves parts~(a)--(c) of Theorem~4'.

\subsection{Part IV: Multi-State Aggregation (Lemma~F)}
\label{sec:lemmaF}

In practice the data contains multiple states, each with its own
$p_{0,s}$, $p_{1,s}$, Chernoff information $\kappa_s$, and per-state
constant $C_s$.  The global $\operatorname{F1}$ score is the expectation
over the state mixture.

\begin{lemma}[Multi-state aggregation]
  \label{lem:aggregation}
  Let there be $S$ states with proportions $\rho_s>0$, $\sum_s\rho_s=1$.
  Within state $s$, the per-state $\operatorname{F1}_s(M)$ satisfies
  \[
  1-\operatorname{F1}_s(M)
  \sim \frac{C_s}{\eta}\,
       \frac{e^{-M\kappa_s}}{\sqrt{2\pi M}}.
  \]
  Then:
  \begin{enumerate}[(i)]
    \item \textbf{Additivity.}
      $\displaystyle
      1-\operatorname{F1}_{\text{global}}(M)
      = \sum_{s=1}^S \rho_s\,(1-\operatorname{F1}_s(M))$.
    \item \textbf{Bottleneck rate.}
      $\displaystyle
      \kappa_{\text{global}} = \min_{1\le s\le S} \kappa_s$.
    \item \textbf{Dominant-state constant.}
      Let $\mathcal{S}_{\min}=\{s:\kappa_s=\kappa_{\text{global}}\}$.  Then
      \[
      \liminf_{M\to\infty}
        e^{M\kappa_{\text{global}}}\sqrt{2\pi M}\,
        \bigl(1-\operatorname{F1}_{\mathcal{A}}(M)\bigr)
      \ge \frac1\eta\sum_{s\in\mathcal{S}_{\min}} \rho_s C_s.
      \]
      States with $\kappa_s>\kappa_{\text{global}}$ contribute only
      $o(e^{-M\kappa_{\text{global}}}/\sqrt{M})$ and are asymptotically
      negligible.
  \end{enumerate}
\end{lemma}

\begin{proof}
  Part~(i) follows from linearity of expectation.  For part~(ii), each state
  contributes $e^{-M\kappa_s}/\sqrt{2\pi M}$.  The slowest exponential decay
  (smallest $\kappa_s$) dominates as $M\to\infty$.  Part~(iii) follows from
  Lemma~E applied per state and summing with weights $\rho_s$; states with
  $\kappa_s>\kappa_{\text{global}}$ are suppressed by
  $e^{-M(\kappa_s-\kappa_{\text{global}})}\to0$.
\end{proof}

\subsection{Numerical Verification}
\label{sec:thm4_numerical}

We verify the theoretical constants numerically across five parameter sets
spanning a range of noise rates and distribution separations.

\subsubsection{Five Parameter Sets}

The following table reports for each $(p_0,p_1,\eta)$ combination: the
Chernoff information $\kappa$, the ratio $2(p_1-p_0)^2/\kappa$ (showing
how much looser the Hoeffding bound is), the optimal constant $C_{\min}$,
the normalised constant $C_{\min}/\eta$, the suboptimality ratio of the
non-adaptive threshold $\theta^*$, and whether the adaptive threshold
matches the lower bound.

\[
\begin{array}{c|c|c|c|c|c|c|c|c|c}
\text{Case} & p_0 & p_1 & \eta & \kappa &
\frac{2(p_1-p_0)^2}{\kappa}
& C_{\min} & \frac{C_{\min}}{\eta}
& \frac{\text{Non-ad.}}{C_{\min}/\eta}
& \text{Adapt match?} \\ \hline
1 & 0.10 & 0.60 & 0.10 & 0.1696 & 2.95
  & 0.4591 & 4.5915 & 1.70 & \text{YES} \\[4pt]
2 & 0.20 & 0.50 & 0.30 & 0.0528 & 3.41
  & 1.3768 & 4.5894 & 1.09 & \text{YES} \\[4pt]
3 & 0.05 & 0.80 & 0.05 & 0.4574 & 2.46
  & 0.1843 & 3.6864 & 2.41 & \text{YES} \\[4pt]
4 & 0.10 & 0.60 & 0.50 & 0.1696 & 2.95
  & 0.8348 & 1.6697 & 1.00 & \text{YES} \\[4pt]
5 & 0.10 & 0.60 & 0.90 & 0.1696 & 2.95
  & 0.5465 & 0.6072 & 1.62 & \text{YES}
\end{array}
\]

\subsubsection{Key Findings}

\begin{enumerate}
  \item \textbf{Adaptive limit matches $C_{\min}/\eta$ exactly.}
    For all five cases, the difference between the adaptive threshold
    constant and the lower bound constant is below $10^{-15}$ (machine
    precision), confirming Theorem~4'(c).
  \item \textbf{Non-adaptive threshold is suboptimal.}
    Using the naive Chernoff point $\theta^*$ (which ignores $\eta$)
    yields constants $1.00$--$2.41\times$ larger than $C_{\min}/\eta$.
    The worst ratio (2.41) occurs for rare noise ($\eta=0.05$) with
    widely separated distributions.
  \item \textbf{Chernoff rate is slower than Hoeffding bound.}
    The ratio $2\Delta^2/\kappa$ ranges from 2.46 to 3.41, confirming
    that $\kappa<2\Delta^2$ for all tested parameters.  The Hoeffding
    exponential bound is looser by a factor of 2.5--3.4, justifying the
    use of the exact Chernoff information.
  \item \textbf{Symmetric case $\eta=0.5$.}
    When $\eta=0.5$, $\theta^\dagger=\theta^*$ and the adaptive and
    non-adaptive thresholds coincide (ratio 1.00).  This is expected
    because $((1-\eta)/\eta)^s=1$ for $\eta=0.5$.
\end{enumerate}

\subsubsection{Finite-$M$ Convergence}

The asymptotic limits are derived as $M\to\infty$.  For finite $M$, the
Bahadur-Rao theorem provides an $O(1/M)$ correction term.  Lemma~B gives
explicit bounds:
\[
1-\operatorname{F1}_{\text{SCX}}(\theta^\dagger)
\le \frac{C_{\min}}{\eta}\,
    \frac{e^{-M\kappa}}{\sqrt{2\pi M}}\,
    \bigl(1 + K_1 M^{-1/2}\bigr),
\]
\[
1-\operatorname{F1}_{\text{SCX}}(\theta^\dagger)
\ge \frac{C_{\min}}{\eta}\,
    \frac{e^{-M\kappa}}{\sqrt{2\pi M}}\,
    \bigl(1 - K_2 M^{-1/2}\bigr),
\]
where $K_1,K_2$ depend only on $p_0,p_1$.  Numerical experiments
(Supplementary Material~S8) confirm that the finite-$M$ F1 score approaches
the asymptotic limit at the predicted $O(1/\sqrt{M})$ rate.

\subsection*{References for Section~S4}

\begingroup
\small
\begin{thebibliography}{20}
\bibitem{bahadur1960deviations} Bahadur, R. R. \& Rao, R. R. On deviations of the sample mean. \textit{Annals of Mathematical Statistics} \textbf{31}(4), 1015--1027 (1960).
\end{thebibliography}
\endgroup

\endinput
